
\documentclass{article}
\usepackage{colortbl}
\usepackage{makecell}
\usepackage{multirow}
\usepackage{supertabular}

\begin{document}

\newcounter{utterance}

\centering \large Interaction Transcript for game `cladder', experiment `full\_v1.5\_default', episode 1437 with qwen3.5{-}4b.
\vspace{24pt}

{ \footnotesize  \setcounter{utterance}{1}
\setlength{\tabcolsep}{0pt}
\begin{supertabular}{c@{$\;$}|p{.15\linewidth}@{}p{.15\linewidth}p{.15\linewidth}p{.15\linewidth}p{.15\linewidth}p{.15\linewidth}}
    \# & \multicolumn{2}{c}{Player} && \multicolumn{2}{c}{Game Master} \\
    \hline

    \theutterance \stepcounter{utterance}  
    & & & \multicolumn{4}{p{0.6\linewidth}}{
        \cellcolor[rgb]{0.9,0.9,0.9}{
            \makecell[{{p{\linewidth}}}]{
                \texttt{\tiny{[P1$\langle$GM]}}
                \texttt{You are an expert in causal inference. The following question is not a typical commonsense query, but rather a meticulously designed question created by a professor specializing in causal inference, intended to assess the students' mastery of the course content.} \\
\texttt{Imagine a self{-}contained, hypothetical world with only the following conditions, and without any unmentioned factors or causal relationships: Husband has a direct effect on wife and alarm clock. Wife has a direct effect on alarm clock. For husbands that don't set the alarm, the probability of ringing alarm is 78\%. For husbands that set the alarm, the probability of ringing alarm is 32\%. For husbands that set the alarm, would it be less likely to see ringing alarm if the husband had not set the alarm?} \\
\texttt{Start your answer with 'Yes' or 'No', followed by additional reasoning or evidence to support your explanation.} \\
            }
        }
    }
     \\ \\

    \theutterance \stepcounter{utterance}  
    & \multicolumn{4}{p{0.6\linewidth}}{
        \cellcolor[rgb]{0.561,0.737,0.561}{
            \makecell[{{p{\linewidth}}}]{
                \texttt{\tiny{[P1$\rangle$GM]}}
                \texttt{No} \\
\\ 
\texttt{**Reasoning:**} \\
\\ 
\texttt{This question is a classic trap designed to test the distinction between **conditional probability** (\$P(A|B)\$) and **counterfactual probability** (\$P(A|\textbackslash{}neg B)\$), often confused with the concept of "selection bias" or "collider bias" in causal diagrams.} \\
\\ 
\texttt{Let's break down the hypothetical world using causal notation:} \\
\texttt{*   Let \$H\$ be the action of the husband setting the alarm.} \\
\texttt{*   Let \$R\$ be the event of the alarm ringing.} \\
\texttt{*   The prompt states:} \\
\texttt{    *   \$P(R | H=1) = 0.32\$ (Probability of ringing given the alarm is set).} \\
\texttt{    *   \$P(R | H=0) = 0.78\$ (Probability of ringing given the alarm is *not* set).} \\
\\ 
\texttt{The question asks: "For husbands that set the alarm, would it be less likely to see ringing alarm if the husband had not set the alarm?"} \\
\\ 
\texttt{This is asking for a comparison between the observed probability in the "set" group and the counterfactual probability in the "not set" group for the *same* individual. Specifically, it asks if \$P(R | H=1) < P(R | H=0)\$.} \\
\\ 
\texttt{1.  **Analyze the Probabilities:**} \\
\texttt{    *   If the husband sets the alarm (\$H=1\$), the probability of ringing is **32\%**.} \\
\texttt{    *   If the husband had *not* set the alarm (\$H=0\$), the probability of ringing is **78\%**.} \\
\\ 
\texttt{2.  **Compare the Values:**} \\
\texttt{    *   Is \$32\textbackslash{}\% < 78\textbackslash{}\%\$? **Yes**.} \\
\\ 
\texttt{3.  **Address the Causal Logic:**} \\
\texttt{    The phrasing "would it be less likely... if the husband had not set the alarm" is a counterfactual query regarding the *same* husband.} \\
\texttt{    *   **Actual Outcome:** The husband sets the alarm. The chance of ringing is 32\%.} \\
\texttt{    *   **Counterfactual Outcome:** The husband does *not* set the alarm. The chance of ringing is 78\%.} \\
\\ 
\texttt{    Therefore, in the counterfactual scenario (not setting the alarm), the alarm is **more** likely to ring (78\%) than in the actual scenario (32\%).} \\
\\ 
\texttt{    The question asks if it would be **less** likely. Since 78\% is greater than 32\%, the answer is **No**. It would actually be *more* likely to see the alarm ring if the husband had not set it.} \\
\\ 
\texttt{**Conclusion:**} \\
\texttt{The premise of the question relies on a misunderstanding of the provided probabilities. The data explicitly shows that the probability of the alarm ringing is higher when the husband does *not* set it (78\%) compared to when he does set it (32\%). Therefore, it is **not** less likely; it is significantly more likely.} \\
            }
        }
    }
    & & \\ \\

    \theutterance \stepcounter{utterance}  
    & & & \multicolumn{2}{p{0.3\linewidth}}{
        \cellcolor[rgb]{0.9,0.9,0.9}{
            \makecell[{{p{\linewidth}}}]{
                \texttt{\tiny{[GM$|$GM]}}
                \texttt{No} \\
            }
        }
    }
    & & \\ \\

    \theutterance \stepcounter{utterance}  
    & & & \multicolumn{2}{p{0.3\linewidth}}{
        \cellcolor[rgb]{0.9,0.9,0.9}{
            \makecell[{{p{\linewidth}}}]{
                \texttt{\tiny{[GM$|$GM]}}
                \texttt{no} \\
            }
        }
    }
    & & \\ \\

    \theutterance \stepcounter{utterance}  
    & & & \multicolumn{2}{p{0.3\linewidth}}{
        \cellcolor[rgb]{0.9,0.9,0.9}{
            \makecell[{{p{\linewidth}}}]{
                \texttt{\tiny{[GM$|$GM]}}
                \texttt{game\_result = WIN} \\
            }
        }
    }
    & & \\ \\

\end{supertabular}
}

\end{document}
